\documentclass[a4paper,12pt,titlepage]{article}

\usepackage{cmap}
\usepackage[T2A]{fontenc}
\usepackage[english,russian]{babel}
\usepackage{geometry,indentfirst,graphicx}
\usepackage[hidelinks]{hyperref}

\usepackage{amssymb,amsmath}
\usepackage{multirow}

\setcounter{secnumdepth}{0}
\linespread{1.5}
\geometry{a4paper,left=30mm,top=20mm,bottom=20mm,right=15mm}

\begin{document}

\hypersetup{pageanchor=false}
\begin{titlepage}
    \begin{center}
        \includegraphics{logo.pdf} \\
        \textsc{\small Московский государственный университет имени М.~В.~Ломоносова \\
        Факультет вычислительной математики и кибернетики \\
        Кафедра математической кибернетики \\}
        \vfill
        \large{Королёв Фёдор Иванович} \\
        \vspace{1cm}
        \textbf{\large Программная реализация поиска минимальных обобщенных полиномов в векторном пространстве} \\
        \vfill
        \textsc{\large Выпускная квалификационная работа} \\
    \end{center}
    \begin{flushright}
        \vfill
        \textbf{Научный руководитель:} \\
        к.ф.-м.н. Бухман~А.~В.
    \end{flushright}
    \begin{center}
        \vfill
        {\small Москва \\
        2022}
    \end{center}
\end{titlepage}

\hypersetup{pageanchor=true}
\tableofcontents
\newpage

\section{Основные понятия}

Пусть $ E_2 $ --- множество $ \{ 0, 1 \} $.\\
Отображение $ f:\ E_2^n \rightarrow E_2 $ назовем булевой функцией от $ n $ переменных.\\
$ \mathbb{P}_2(n) $ --- множество всех булевых функций от $ n $ переменных.\\
Если $ \sigma \in E_2 $, то введем обозначение
\begin{equation*}
    x^\sigma =
    \begin{cases}
        x & \text{при $ \sigma = 1 $}\\
        \overline x & \text{при $ \sigma = 0 $,}
    \end{cases}
\end{equation*}
выражение $ x^\sigma $ будем называть \textit{литералом}.\\
\textit{Элементарная конъюнкция} ранга $ r $ --- это выражение вида $ x_{i_1}^{\sigma_1} \dots x_{i_r}^{\sigma_r} $, где $ x_{i_1}, \dots, x_{i_r} $ --- различные переменные и $ \sigma_1, \dots, \sigma_r \in E_2 $. Элементарной конъюнкцией ранга 0 назовем константу 1.\\
\textit{Обобщенный полином} --- выражение вида $ \bigoplus\limits_{i = 1}^s K_i $, где $ K_i $ --- элементарные конъюнкции.\\
\textit{Длиной обобщенного полинома} назовем количество элементарных конъюнкций в нем.\\
$ \mathcal{P}(n) $ --- множество всех обобщенных полиномов от $ n $ переменных.

\section{Введение}

\subsection{Задача минимизации обобщенных полиномов}

Одним из возможных способов задать булеву функцию $ f $ от $ n $ переменных является задание ее таблицы истинности, количество строк которой равно $ 2^n $. Отсюда следует, что мощность класса булевых функций от $ n $ переменных $ \left| \mathbb{P}_2(n) \right| $ равна $ 2^{2^n} $, поскольку на каждую строку таблицы истинности булева функция может ответить нулем или единицей, то есть элементами из множества $ \{ 0,\ 1 \} $ мощности 2.

Элементы класса обобщенных полиномов $ \mathcal{P}(n) $ от $ n $ переменных задаются, в соответствии с представленным определением, исключающим или, примененным ко всем вхождениям в полином элементарных конъюнкций. Мощность класса элементарных конъюнкций равна $ 3^n $, поскольку $ n $ переменных могут принять одно из значений множества $ \{ 1, x_i, \overline x_i \},\ i = \overline{1,\ n} $. Тогда каждый конъюнкт из класса элементарных конъюнкций может входить или не входить в обобщенный полином, поэтому мощность множества всех обобщенных полиномов $ n $ переменных $ \left| \mathcal{P}(n) \right| $ равна $ 2^{3^n} $.

Поскольку $ \left| \mathbb{P}(n) \right| \ne \left| \mathcal{P}(n) \right| $, то между классом булевых функций и классом обобщенных полиномов не существует взаимнооднозначного соответствия. В таком случае введем отношение эквивалентности в классе обобщенных полиномов. Назовем два обобщенных полинома $ p_1,\ p_2 \in \mathcal{P}(n) $ эквивалентными, если они реализуют одну булеву функцию.

Это отношение порождает классы эквивалентности в пространстве обобщенных полиномов. Сформулируем математическую задачу: поиск минимального (в смысле длины) полинома в классе эквивалентности, порожденном элементом $ p $, где $ p $ --- входной обобщенный полином данной задачи. Его представлением будем считать набор элементарных конъюнкций, входящих в него. Выходными данными поставленной задачи будем считать минимальный представитель класса эквивалентности, порожденного входным полиномом $ p $.

Рассмотрим некоторые характеристики представления булевых функций обобщенными полиномами.\\
Для $ P \in \mathcal{P}(n) $ обозначим длину полинома $ l(P) $. Введем функционал
$$ l(f) = \min\limits_{P_f \in \mathcal{P}_f} l(P_f), $$
обозначающий минимальную длину обобщенного полинома среди всех, реализующих булеву функцию $ f $. Функцией Шеннона сложности в классе обобщенных полиномов называется функция
$$ L_\text{о.п.} = \max\limits_{f \in \mathbb{P}_2(n)} l(f). $$

Существуют верхняя и нижняя оценки для функции Шеннона в классе обобщенных полиномов, представим их \cite{selezn}:
$$ \frac{2^n}{n \log_2 3} \le L_\text{о.п.}(n) \le 2 \cdot \frac{2^n}{n} (1 + \ln n) $$

\subsection{Развитие исследований}

Интерес задачи минимизации обобщенных полиномов представляют преимущества, которыми они обладают в сравнении с дизъюнктивными нормальными формами. В производстве цифровых схем на базе логических вентилей \textsc{AND} и \textsc{XOR} уменьшены площадь схемы и время задержки \cite{delay} в сравнении со схемами на базе \textsc{AND} и \textsc{OR}. Текущее применение \textsc{AND-XOR} выражений заключается в синтезе таких двухуровневых схем, как обратимые логические\footnote{схемы, использующие вентили с характерным низким энергопотреблением, реализующие обратимые вычисления \cite{revsynth}} \cite{reversible} и квантовые \cite{quantum}.

Некоторые исследователи предлагают алгоритмы, находящие точный полиномиальный минимум \cite{min-tau,exact6,exact} функции, однако до сих пор не был найден эффективный метод решения таким образом поставленной задачи для булевых функций более восьми переменных в общем случае \cite{exact8}, поэтому особый интерес представляет эвристический подход к минимизации обобщенных полиномов.

Научным сообществом было предложено множество эвристических алгоритмов минимизации, не гарантирующих минимальности полученного в качестве результата полинома, однако позволяющих добиться высокого быстродействия в сравнении с подходами с поиском точного минимума. Примерами могут служить алгоритмы Сасао \cite{exmin2}, Стергиу-Вудурис-Папаконстантину \cite{mvesopmin}, Мищенко-Перковски \cite{exorcism4}. Стоит упомянуть, что некоторые алгоритмы (например \textsc{EXMIN2} \cite{exmin2}) принимают на вход дизъюнктивную нормальную форму, но на первом алгоритмическом этапе занимаются приведением ее к обобщенному виду.

Существуют и алгоритмы минимизации \cite{exorcism-mv3}, удовлетворяющие условию крайне высокого быстродействия. Такой подход находит применение в алгоритмах, использующих минимизацию обобщенных полиномов в цикле другой минимизирующей программы, например в машинном обучении \cite{machine-learning}.

\section{Постановка задачи}

Зафиксируем следующие задачи:
\begin{enumerate}
    \item Предложить эвристический подход к решению задачи минимизации обобщенных полиномов
    \item Написать программу, реализующую предложенный алгоритм
    \item Привести оценку качества ответа и быстродействия реализованной программы
    \item Сравнить предложенный алгоритм с известными алгоритмами минимизации обобщенных полиномов
\end{enumerate}

\section{Основная часть}

\subsection{Метод прибавления нулевых полиномов}

Задан исходный обобщенный полином в виде $ \bigoplus\limits_i K_i $. Генерируется некоторое множество нулевых полиномов --- обобщенных полиномов, тождественно равных нулю. Поскольку $ a \oplus 0 \equiv a $, то операция прибавления элементов такого множества к исходному полиному будет сохранять класс эквивалентности, при этом может повлечь за собой сокращение длины полинома (поскольку для некоторых конъюнктов может сработать правило сокращения $ a \oplus a \equiv 0 $).

\subsubsection{Описание начального множества нулевых полиномов}
Под начальным множеством нулевых полиномов здесь и далее будем понимать множество нулевых полиномов, сгенерированных следующим образом.
Представим себе элементарную конъюнкцию из $ n $ множителей. На позиции i-того литерала выставим один из четырех элементов множества $ \{ 0, 1, x_i, \overline{x_i} \} $. Очевидно, что конъюнкция тождественно равна 0 тогда и только тогда, когда на месте хотя бы одного литерала выставлен 0. Пусть на места всех $ n $ литералов выставлены элементы описанного множества, при этом хотя бы один из них равен $ 0 \equiv 1 \oplus x_i \oplus \overline x_i $, тогда в силу дистрибутивности конъюнкции относительно исключающего или можно раскрыть скобки и получить выражение, являющееся, в силу тождественного равенства нулю, нулевым обобщенным полиномом. Все начальное множество нулевых полиномов, построенное на множестве $ \{ 0, 1, x_i, \overline{x_i} \} $, будет состоять из $ 4^n $ (всех вариантов расстановки элементов множества на $ n $ позиций) минус $ 3^n $ (всех вариантов расстановки элементов $ \{ 1, x_i, \overline{x_i} \} $ на $ n $ позиций, поскольку выставив исключительно элементы данного подмножества на места литералов выражение не будет удовлетворять необходимому условию тождественного равенства нулю). Итого, мощность начального множества $ \left| B \right| = 4^n - 3^n. $

\subsubsection{Алгоритм минимизации}

В настоящей работе предложен следующий эвристический алгоритм минимизации обобщенных полиномов:
\begin{enumerate}
    \item Генерируем описанное выше начальное множество нулевых обобщенных полиномов.
    \item Расширим множество $ C_{\left| B \right|}^2 $ комбинациями всевозможных пар его элементов $ B $, $ C_{\left| B \right|}^3 $ тройками по аналогии, и так далее до разумного с точки зрения потребления оперативной памяти предела.
    \item Запускаем поиск сокращенного эквивалентного полинома:
    \begin{enumerate}
        \item \label{loop:begin} Рассматриваем результаты прибавлений всех нулевых полиномов из расширенного сочетаниями множества, имеющих с текущим представлением сокращаемого полинома хотя бы один общий конъюнкт.
        \item \label{loop:end} Выбираем из всех вариантов тот, который позволил сократить длину полинома на наибольшее значение. В случае одинаковых сокращений выбирается сокращение с последним вошедшим в множество нулевым полиномом.
        \item Алгоритм заканчивает работу, если на этапе \ref{loop:end} не нашлось ни единого кандидата, иначе возвращается к пункту \ref{loop:begin}.
    \end{enumerate}
\end{enumerate}

Алгоритм возвращает сокращенный обобщенный полином в виде $ \bigoplus\limits_i K_i $.

Настоящий алгоритм отличается от его предыдущей версии тем, что за счет предварительного создания сочетаний элементов начального множества удалось значительно улучшить скорость работы программы, что позволило пробовать использовать более <<глубокие>> последовательности прибавлений (ранее прибавлялись только сами элементы множества и их всевозможные пары).

\subsubsection{Реализация программы}

Была написана программа на языке \texttt{C++}, исходный код которой доступен по ссылке: \url{https://github.com/skajida/diploma}.

\subsection{Тестирование}

Как упоминалось выше, основным фактором, мешающим дополнять множество нулевых полиномов комбинациями элементов начального множества (пусть далее его мощность $ n $), является ограниченность оперативной памяти. При очередной попытке расширить множество $ C_n^k $ комбинациями из $ k $ элементов начального множества, программа оценивает ожидаемый запрос памяти и, если он превышает \texttt{8GB}, то отклоняет попытку расширения.

Таким образом удалось расширить множество в пространстве размерности 3 до комбинаций из 5 элементов начального множества, а в пространстве размерности 4 до комбинаций из 3 элементов. Для размерности 5 по оценке потребовалось бы порядка \texttt{80GB} оперативной памяти для хранения самого множества, комбинаций из 2 и 3 его элементов, поэтому для размерностей выше 4 проэкспериментировать с расширением до всевозможных троек не удалось.

По этой причине в настоящей работе будут обсуждаться результаты работы разных подходов именно для размерностей пространств 3 и 4.

В статье Папаконстантину \cite{513100} был озвучен следующий тезис: <<была написана программа, которой на вход подавались 100 случайно сгенерированных функций 5-ти переменных со средним количеством конъюнктов равным 13>>, <<средняя найденная алгоритмом длина равна 5.9, что лучше, чем результат около 6.4 у \cite{beat1, beat2}>>.

Так как для размерности 5 настоящему алгоритму не удается расширить множество нулевых полиномов до троек, то используем данный подход к тестированию работы программы для размерностей 3 и 4 с расширением сочетаниями до пятерок и троек соответственно. В таблице \ref{table_results} отображены длины входных и выходных данных программы в тестах со 100 случайно сгенерированными полиномами средней длины 7, 10, и 13 соответственно. Длина отдельно взятого полинома из входных данных выбиралась из нормального распределения с математическим ожиданием, равным желаемой средней длине, и дисперсией, равной 1. Конъюнкты подбирались из равномерного распределения по их лексикографическому индексу до тех пор, пока длина полинома не будет равной заявленной.

\begin{table}[h!]
\centering
\begin{tabular}{ |l||r|r||r|r||r|r||r|r||r|r||r|r| } \hline
\multirow{4}{*}{\bf длина} & \multicolumn{12}{c|}{\bf количество полиномов} \\ \cline{2-13}
    & \multicolumn{2}{c||}{3-dim} & \multicolumn{2}{c||}{4-dim} & \multicolumn{2}{c||}{3-dim} & \multicolumn{2}{c||}{4-dim} & \multicolumn{2}{c||}{3-dim} & \multicolumn{2}{c|}{4-dim} \\ \cline{2-13}
    & \texttt{in} & \texttt{out} & \texttt{in} & \texttt{out} & \texttt{in} & \texttt{out} & \texttt{in} & \texttt{out} & \texttt{in} & \texttt{out} & \texttt{in} & \texttt{out} \\ \hline \hline
1   & 0     & 9     & 0     & 0     & 0     & 7     & 0     & 0     & 0     & 10    & 0     & 0     \\ \hline
2   & 0     & 61    & 0     & 4     & 0     & 75    & 0     & 4     & 0     & 63    & 0     & 4     \\ \hline
3   & 0     & 30    & 0     & 31    & 0     & 18    & 0     & 23    & 0     & 27    & 0     & 24    \\ \hline
4   & 3     & 0     & 0     & 40    & 0     & 0     & 0     & 41    & 0     & 0     & 0     & 42    \\ \hline
5   & 5     & 0     & 3     & 22    & 0     & 0     & 0     & 22    & 0     & 0     & 0     & 21    \\ \hline
6   & 26    & 0     & 22    & 3     & 0     & 0     & 0     & 10    & 0     & 0     & 0     & 8     \\ \hline
7   & 43    & 0     & 39    & 0     & 2     & 0     & 0     & 0     & 0     & 0     & 0     & 1     \\ \hline
8   & 17    & 0     & 31    & 0     & 4     & 0     & 9     & 0     & 0     & 0     & 0     & 0     \\ \hline
9   & 5     & 0     & 5     & 0     & 23    & 0     & 31    & 0     & 0     & 0     & 0     & 0     \\ \hline
10  & 1     & 0     & 0     & 0     & 41    & 0     & 35    & 0     & 1     & 0     & 2     & 0     \\ \hline
11  & 0     & 0     & 0     & 0     & 25    & 0     & 18    & 0     & 4     & 0     & 4     & 0     \\ \hline
12  & 0     & 0     & 0     & 0     & 4     & 0     & 7     & 0     & 24    & 0     & 28    & 0     \\ \hline
13  & 0     & 0     & 0     & 0     & 1     & 0     & 0     & 0     & 38    & 0     & 33    & 0     \\ \hline
14  & 0     & 0     & 0     & 0     & 0     & 0     & 0     & 0     & 30    & 0     & 28    & 0     \\ \hline
15  & 0     & 0     & 0     & 0     & 0     & 0     & 0     & 0     & 3     & 0     & 5     & 0     \\ \hline \hline
\textbf{mean len}   & 6.85  & 2.21  & 7.13  & 3.89  & 9.99  & 2.11  & 9.83  & 4.11  & 13.01 & 2.17  & 12.96 & 4.08  \\ \hline
\textbf{reduct}  & \multicolumn{2}{c||}{\tt -68\%}  & \multicolumn{2}{c||}{\tt -45\%} & \multicolumn{2}{c||}{\tt -79\%}  & \multicolumn{2}{c||}{\tt -58\%}   & \multicolumn{2}{c||}{\tt -83\%}  & \multicolumn{2}{c|}{\tt -69\%} \\ \hline
\textbf{time}       & \multicolumn{2}{c||}{\tt 4m 44s}  & \multicolumn{2}{c||}{\tt 15m 47s} & \multicolumn{2}{c||}{\tt 8m 25s}  & \multicolumn{2}{c||}{\tt 24m 09s}  & \multicolumn{2}{c||}{\tt 8m 43s}  & \multicolumn{2}{c|}{\tt 30m 38s}  \\ \hline
\end{tabular}
\caption{Результаты работы программы на полиномах длины 7, 10, 13}
\label{table_results}
\end{table}

\subsection{Сравнение с предыдущей реализацией}

Предыдущая версия программы работала только с начальным множеством и его всевозможными парами, качественное сравнение результатов представлено в таблице \ref{table_comparison}, где под <<глубиной>> (depth) подразумевается максимальное $ k $ рассматриваемых $ C_n^k $ комбинаций элементов множества.

\begin{table}[h!]
\centering
\begin{tabular}{ |l||r|r||r|r|} \hline
\multirow{3}{*}{\bf длина входа} & \multicolumn{4}{c|}{\bf средняя длина выхода/сокращение} \\ \cline{2-5}
    & \multicolumn{2}{c||}{3-dim} & \multicolumn{2}{c|}{4-dim} \\ \cline{2-5}
    & \texttt{2-depth} & \texttt{5-depth} & \texttt{2-depth} & \texttt{3-depth} \\ \hline \hline
7   & 2.95/-57\%  & 2.21/-68\%  & 4.98/-30\%  & 3.89/-45\%  \\ \hline
10  & 2.93/-71\%  & 2.11/-79\%  & 6.11/-38\%  & 4.11/-58\%  \\ \hline
13  & 2.97/-77\%  & 2.17/-83\%  & 6.24/-52\%  & 4.08/-69\%  \\ \hline \hline
\textbf{ср.преимущество}    & \multicolumn{2}{c||}{12\%}  & \multicolumn{2}{c|}{43\%}   \\ \hline
\end{tabular}
\caption{Сравнение с малой глубиной}
\label{table_comparison}
\end{table}

\subsection{Сравнение с опубликованными эвристическими алгоритмами}

Поскольку для размерности выше 4 алгоритм не использует глубину больше двух, то с точки зрения качества ответа он не претендует на конкуренцию с известными эвристическими алгоритмами минимизации \cite{exmin2, exorcism4, exorcism-mv3, grmin2}.

Также в 1992 году была опубликована работа Сасао \cite{4var-min}, в которой описаны свойства минимальных полиномов размерности 4, откуда следует, что все булевы функции представимы обобщенными полиномами длины строго меньше 7. Однако, в таблице \ref{table_results} видно, что из 300 случайно сгенерированных полиномов один из них был приведен к полиному длины 7, что доказывает, что предложенный алгоритм не находит точных полиномиальных минимумов.

С точки зрения производительности алгоритм также оставляет желать лучшего: несмотря на то, что с предыдущей версии была проведена значительная работа по увеличению скорости работы, что позволило рассматривать более <<глубокие>> сочетания начального множества, алгоритм на полиномах размерности 4 тратит в среднем по 14.1 секунд (исходя из временных замеров, представленных в таблице \ref{table_results}), что не идет в сравнение с алгоритмами \cite{exmin2, exorcism4, grmin2}, тем более со специально разработанным для встраивания в другие алгоритмы методом \textsc{EXORCISM-MV3} \cite{exorcism-mv3}.

Данные замечания приводят к выводу о том, что предложенный метод нуждается в значительной доработке, чтобы конкурировать с опубликованными эвристическими алгоритмами минимизации обобщенных полиномов.

\section{Полученные результаты}

\begin{enumerate}
    \item Был предложен эвристический подход к минимизации обобщенных полиномов путем прибавления нулевых полиномов
    \item Была написана программа, реализующая предложенный алгоритм
    \item Была проведена оценка качества программы на случайных обобщенных полиномах из пространств обобщенных полиномов размерностей 3 и 4
\end{enumerate}

\begin{raggedright}
    \addcontentsline{toc}{section}{Литература}
    \begin{thebibliography}{99}
        \bibitem{delay} Kalay~U., Hall~D.~V., Perkowski~M.~A. A minimal universal test set for self-test of EXOR-sum-of-products circuits //IEEE Transactions on Computers. – 2000. – Т.~49. – №.~3. – С.~267-276.
        \bibitem{reversible} Yang~G. et al. Majority-based reversible logic gates //Theoretical computer science. – 2005. – Т.~334. – №.~1-3. – С.~259-274.
        \bibitem{quantum} Iwama~K., Kambayashi~Y., Yamashita~S. Transformation rules for designing CNOT-based quantum circuits //Proceedings of the 39th annual Design Automation Conference. – 2002. – С.~419-424.
        \bibitem{revsynth} Shende~V.~V. et al. Reversible logic circuit synthesis //Proceedings of the 2002 IEEE/ACM international conference on Computer-aided design. – 2002. – С.~353-360.
        \bibitem{min-tau} Hirayama~T., Nishitani~Y., Sato~T. A faster algorithm of minimizing AND-EXOR expressions //IEICE transactions on fundamentals of electronics, communications and computer sciences. – 2002. – Т.~85. – №.~12. – С.~2708-2714.
        \bibitem{exact6} Gaidukov~A. Algorihm to derive minimum ESOP for 6-variable function //5th International Workshop on Boolean Problems, Sept.~2002. – 2002.
        \bibitem{exact} Sasao~T. An exact minimization of AND-EXOR expressions using reduced covering functions //Proc. of the Synthesis and Simulation Meeting and International Interchange. – 1993. – С.~374-383.
        \bibitem{exact8} Stergiou~S., Papakonstantinou~G. Exact minimization of ESOP expressions with less than eight product terms //Journal of Circuits, Systems, and Computers. – 2004. – Т.~13. – №.~01. – С.~1-15.
        \bibitem{exmin2} Sasao~T. EXMIN2: A simplification algorithm for exclusive-OR-sum-of-products expressions for multiple-valued-input two-valued-output functions //IEEE Transactions on Computer-Aided Design of Integrated Circuits and Systems. – 1993. – Т.~12. – №.~5. – С.~621-632.
        \bibitem{mvesopmin} Stergiou~S., Voudouris~D., Papakonstantinou G. Multiple-value exclusive-or sum-of-products minimization algorithms //IEICE TRANSACTIONS on Fundamentals of Electronics, Communications and Computer Sciences. – 2004. – Т.~87. – №.~5. – С.~1226-1234.
        \bibitem{exorcism4} Mishchenko~A., Perkowski~M. Fast heuristic minimization of exclusive-sums-of-products. – 2001.
        \bibitem{exorcism-mv3} Song~N., Perkowski~M. New fast approach to approximate ESOP minimization for incompletely specified multi-output functions //Proc. IFIP 10.5 Workshop on Application of Reed-Muller Expansions in Circuit Design. – 1997. – С. 61-72.
        \bibitem{machine-learning} Perkowski~M. et al. Application of ESOP minimization in machine learning and knowledge discovery //Proc. Reed Muller. – 1995. – Т. 95.
        \bibitem{selezn} Селезнева~С.~Н., Шуплецов~М.~С., Дайняк~А.~Б. Булевы функции и полиномы //Москва – 2006.
        \bibitem{513100} Papakonstantinou~G. Minimization of modulo-2 sum of products //IEEE Transactions on Computers. – 1979. – Т. 28. – №. 02. – С. 163-167.
        \bibitem{beat1} Even~S., Kohavi~I., Paz~A. On minimal modulo 2 sums of products for switching functions //IEEE Transactions on Electronic Computers. – 1967. – №. 5. – С. 671-674.
        \bibitem{beat2} Bioul~G. Minimization of ring-sum expansions of Boolean functions //Philips Res. Repts. – 1973. – Т. 28. – С. 17-36.
        \bibitem{grmin2} Debnath~D., Sasao~T. GRMIN2: A heuristic simplification algorithm for generalised Reed-Muller expressions //IEE Proceedings-Computers and Digital Techniques. – 1996. – Т.~143. – №.~6. – С.~376-384.
        \bibitem{4var-min} Koda~N., Sasao~T. Four‐variable AND‐EXOR minimum expressions and their properties //Systems and computers in Japan. – 1992. – Т. 23. – №. 10. – С. 27-41.
    \end{thebibliography}
\end{raggedright}

\end{document}
